\subsection{庞加莱群胚的函子性}

设 X、Y 是拓扑空间，\(f: X \to Y\) 是连续映射。映射 \(c \mapsto f \circ c\) 是连续映射，记为 \(\Lambda(f)\)，从 \(\Lambda(X) = \mathcal{C}_c(\mathbf{I}; X)\) 到 \(\Lambda(Y) = \mathcal{C}_c(\mathbf{I}; Y)\)（I, p.~132, 引理）。通过取子集，它定义下列连续映射

\[
\Lambda_ {x} (f) \colon \Lambda_ {x} (\mathrm{X}) \to \Lambda_ {f (x)} (\mathrm{Y}),
\]

对 \(x \in X\) 连续，

\[
\Lambda_ {x, y} (f) \colon \Lambda_ {x, y} (\mathrm{X}) \to \Lambda_ {f (x), f (y)} (\mathrm{Y}),
\]

对 \(x, y \in \mathrm{X}\) 连续，并且

\[
\Omega (f) \colon \Omega (\mathrm{X}) \to \Omega (\mathrm{Y}).
\]

对 \(x \in X\)，映射 \(\Lambda_{x,x}(f)\) 也记为 \(\Omega_x(f)\)。取道路连通分支后（III, p.~290），由映射 \(\Lambda_{x,y}(f)\) 得到映射

\[
\varpi_ {x, y} (f) \colon \varpi_ {x, y} (\mathrm{X}) \to \varpi_ {f (x), f (y)} (\mathrm{Y}).
\]

设 \(x, y, z\) 是 X 中的点，\(c \in \Lambda_{x,y}(X)\) 和 \(d \in \Lambda_{y,z}(X)\) 是道路；根据道路拼接的定义，有

\[
f \circ (c * d) = (f \circ c) * (f \circ d).
\]

取严格同伦类后得到关系

\[
\varpi_ {x, z} (f) (\gamma \delta) = (\varpi_ {x, y} (f) (\gamma)) (\varpi_ {y, z} (f) (\delta))
\]

 对所有 \(\gamma \in \varpi_{x,y}(\mathrm{X})\) 及每个 \(\delta \in \varpi_{y,z}(\mathrm{X})\) 都成立。因此，连续映射 \(f\) 以及对 \(\varpi_{x,y}(f)\) 和 \(x\) 的映射 \(y \in \mathrm{X}\)，定义了从群胚 \(\varpi(\mathrm{X})\) 到群胚 \(\varpi(\mathrm{Y})\) 的态射（II, p.~161, 定义 3）。记为 \(\varpi(f)\)，称为由连续映射 \(f\) 诱导的庞加莱群胚态射。特别地，若 \(x\) 是 X 中的一点，则映射 \(\varpi_{x,x}(f)\) 是从群 \(\pi_1(\mathrm{X}, x)\) 到群 \(\pi_1(\mathrm{Y}, f(x))\) 的同态；该同态也记为 \(\pi_1(f, x)\)。

注 1. --- 若在集合 \(\varpi_{x,y}(f)\) 和 \(\varpi_{x,y}(\mathrm{X})\) 上赋予 \(\varpi_{f(x),f(y)}(\mathrm{Y})\) 和 \(\mathcal{C}_c(\mathbf{I};\mathrm{X})\) 上紧开拓扑的商拓扑，则同态 \(\mathcal{C}_c(\mathbf{I};\mathrm{Y})\) 是连续的。

为简化记号，当 \(x\)、\(y \in X\) 且 \(\gamma \in \varpi_{x,y}(X)\) 时，我们有时用 \(f_{*}(\gamma)\) 表示 \(\varpi_{x,y}(f)(\gamma)\) 中的元素 \(\varpi_{f(x),f(y)}(Y)\)。

设 X、Y、Z 是拓扑空间，\(f: X \to Y\) 和 \(g: Y \to Z\) 是连续映射。对 X 中的每条道路 c，有 \((g \circ f) \circ c = g \circ (f \circ c)\)。由此有

\[
\varpi (g \circ f) = \varpi (g) \circ \varpi (f).
\]

设 X 和 Y 是拓扑空间，\(f_{0}\) 和 \(f_{1}\) 是从 X 到 Y 的连续映射，\(\sigma\colon X\times I\to Y\) 是连接 \(f_{0}\) 与 \(f_{1}\) 的同伦。由于从 \((c,t)\mapsto c(t)\) 到 X 的映射 \(\mathcal{C}_{c}(\mathbf{I};\mathrm{X})\times\mathbf{I}\)（III, p.~257, 命题 1）连续，由 \(\mathcal{C}_{c}(\mathbf{I};\mathrm{X})\times\mathbf{I}\times\mathbf{I}\) 给出的从 \((c,t,s)\mapsto\sigma(c(t),s)\) 到 Y 的映射也连续。因此，映射 \(\Sigma\colon(c,s)\mapsto\sigma(c(\cdot),s)\) 是从 \(\mathcal{C}_{c}(\mathbf{I};\mathrm{X})\times\mathbf{I}\) 到 \(\mathcal{C}_{c}(\mathbf{I};\mathrm{Y})\) 的连续映射（上引文）。映射 \(\Sigma\) 是连接映射 \(\Lambda(f_{0})\) 与 \(\Lambda(f_{1})\) 的同伦。将其限制到圈空间，映射 \(\Sigma\) 诱导出同伦 \(\Omega(\mathrm{X})\times\mathbf{I}\to\Omega(\mathrm{Y})\)，连接 \(\Omega(f_{0})\) 与 \(\Omega(f_{1})\)。令 x 为 X 中的一点；假设 \(\sigma\) 是在 x 处的带基点同伦，并置 \(y=f_{0}(x)=f_{1}(x)\)。于是映射 \(\Sigma\) 诱导出从 \(\Omega_{x}(\mathrm{X})\times\mathbf{I}\) 到 \(\Omega_{y}(\mathrm{Y})\) 的连续映射，它是以 \(e_{x}\) 为基点的带基点同伦，连接映射 \(\Omega_{x}(f_{0})\) 与 \(\Omega_{x}(f_{1})\)。

命题 3. --- 设 X 和 Y 是拓扑空间，\(f_{0}\) 和 \(f_{1}\) 是从 X 到 Y 的连续映射，\(\sigma\colon\mathbf{X}\times\mathbf{I}\to\mathbf{Y}\) 是连接 \(f_{0}\) 与 \(f_{1}\) 的同伦。令 x 为 X 中的一点；置 \(y_{0}=f_{0}(x)\)、\(y_{1}=f_{1}(x)\)，并将 \(\delta\in\varpi_{y_{0},y_{1}}(\mathbf{Y})\) 记为由 \(d(t)=\sigma(x,t)\)（\(t\in\mathbf{I}\)）定义的道路 d 的类。对所有 \(\gamma\in\pi_{1}(\mathbf{X},x)\)，有 \((f_{1})_{*}(\gamma)=\delta^{-1}\big((f_{0})_{*}(\gamma)\big)\delta\)。

设 c 是 X 中 x 处的一个圈，\(\gamma\) 是它的严格同伦类。置 \(\gamma_{0} = (f_{0})_{*}(\gamma)\) 和 \(\gamma_{1} = (f_{1})_{*}(\gamma)\)；它们分别是 \(f_{0} \circ c\) 和 \(f_{1} \circ c\) 的严格同伦类。对 \((t, s) \in \mathbf{I} \times \mathbf{I}\)，置 \(\varphi(t, s) = \sigma(c(t), s)\)。对所有 \(t \in I\)，有 \(\varphi(t, 0) = (f_{0} \circ c)(t)\)、\(\varphi(t, 1) = (f_{1} \circ c)(t)\) 且 \(\varphi(0, t) = \varphi(1, t) = d(t)\)。因此关系 \(\gamma_{0}\delta = \delta\gamma_{1}\) 由下述引理得到。

引理 1. --- 设 Y 是一个拓扑空间，\(\varphi\colon\mathbf{I}\times\mathbf{I}\to\mathbf{Y}\) 是连续映射。对 \(t\in\mathbf{I}\)，置 \(c_{0}(t)=\varphi(t,0)\)、\(c_{1}(t)=\varphi(t,1)\)、\(d_{0}(t)=\varphi(0,t)\) 和 \(d_{1}(t)=\varphi(1,t)\)。道路 \(c_{0}*d_{1}\) 与 \(d_{0}*c_{1}\) 严格同伦。

令 c 为在 \(I \times I\) 中通过拼接道路 \(t \mapsto (t,0)\) 和 \(t \mapsto (1,t)\) 得到的道路；令 d 为在 \(I \times I\) 中通过拼接道路 \(t \mapsto (0, t)\) 和 \(t \mapsto (t, 1)\) 得到的道路。道路 c 和 d 具有相同的起点 \(c\)、\(d\)、\((0, 0)\) 和相同的终点 \((1, 1)\)。从 \((t, s) \mapsto (1 - s)c(t) + sd(t)\) 到 \(\mathbf{I} \times \mathbf{I}\) 的映射 \(\mathbf{I} \times \mathbf{I}\) 是连接 c 与 d 的严格同伦。有 \(c\)、\(d\)、\(c_0 * d_1 = \varphi \circ c\) 和 \(d_0 * c_1 = \varphi \circ d\)；因此这两条道路严格同伦。

推论 1. --- 设 X 和 Y 是拓扑空间，x 是 X 中的一点，\(f_0\) 和 \(f_1\) 是从 X 到 Y 的连续映射。若存在在 x 处连接 \(f_0\) 与 \(f_1\) 的带基点同伦，则有 \(\pi_1(f_0, x) = \pi_1(f_1, x)\)。

设 \(\sigma\) 是在 \(x\) 处连接 \(f_0\) 与 \(f_1\) 的带基点同伦。用命题 3 的记号，\(\delta\) 是像为 \(f_0(x) = f_1(x)\) 的常值道路的类。断言由此得到。

推论 2. --- 设 X 和 Y 是拓扑空间，\(f\colon X \to Y\) 是同伦等价。对 X 中的每个点 \(x\)，同态

\[
\pi_ {1} (f, x) \colon \pi_ {1} (\mathrm{X}, x) \to \pi_ {1} (\mathrm{Y}, f (x))
\]

是同构。

设 \(g\)、\(f\) 是从 Y 到 X 的连续映射，是 f 的同伦逆。令 \(x\) 为 X 中的一点。将命题 3 应用于同伦的映射 \(\mathrm{Id}_{\mathrm{X}}\) 和 \(g \circ f: \mathrm{X} \to \mathrm{X}\)，可知映射 \(\pi_1(g \circ f, x)\) 是从群 \(\pi_1(\mathrm{X}, x)\) 到群 \(\pi_1(\mathrm{X}, g \circ f(x))\) 的同构。由于

\[
\pi_ {1} (g \circ f, x) = \pi_ {1} (g, f (x)) \circ \pi_ {1} (f, x),
\]

同态 \(\pi_{1}(f,x)\) 是单射，而同态 \(\pi_{1}(g,f(x))\) 是满射。由于映射 g 也是同伦等价，同态 \(\pi_{1}(g,f(x))\) 是单射；因而它是同构。因此，\(\pi_{1}(f,x)\) 是同构。

例. --- 设 G 是一个拓扑群，e 是它的单位元。对 G 中的每个点 g，左平移和右平移 \(x \mapsto gx\) 及 \(x \mapsto xg\) 都是 G 到自身的同胚（TG, III, p.~2），并将 e 映为 g。由推论 2，它们诱导从 \(\pi_{1}(G, e)\) 到 \(\pi_{1}(G, g)\) 的同构。这些同构不一定相等（IV, p.~459, 习题 1）。

推论 3. --- 设 X 是一个同伦等价于一点的拓扑空间。对 X 中的每个点 x，群 \(\pi_{1}(X,x)\) 退化为单位元。

推论 3 特别适用于 X 是 n 维数值空间 \(R^{n}\) 的情形，更一般地，适用于 X 是 \(R^{n}\) 的一个关于其某个点呈星形的子集的情形（III, p.~234）。

命题 4. --- 设 X 是拓扑空间族 \((\mathrm{X}_{j})_{j\in\mathrm{J}}\) 的积空间。由态射族 \(\varpi(\mathrm{X})\) 定义的从群胚 \(\varpi(\mathrm{X}_{j})\) 到群胚 \(j\in J\)（\((\varpi(\mathrm{pr}_{j}))_{j\in\mathrm{J}}\)）的积的态射是同构。

将此群胚态射记为 \(\varphi\)。它取顶点后诱导的映射是恒等映射 \(X \to \prod_{j} X_{j}\)。设 \(x = (x_{j})\) 和 \(y = (y_{j})\) 是 X 中的两个点。记 \(X\)、\(\varphi_{x,y}\) 为由 \(\varpi_{x,y}(X)\) 诱导的从 \(\prod_{j} \varpi_{x_{j},y_{j}}(X_{j})\) 到 \(\varphi\) 的映射。若对每个 \(j \in J\)，\(c_{j}: I \to X_{j}\) 是连接 \(x_{j}\) 与 \(y_{j}\) 的道路，则映射 \(t \mapsto (c_{j}(t))\) 是 X 中连接 x 与 y 的道路（TG, I, p.~25, 命题 1）。这证明了 \(X\)、\(x\)、\(y\)、\(\varphi_{x,y}\)\(c\)、\(d\)、\(X\)、\(x\)、\(y\)、 是满射。设 c 和 d 是 X 中连接 x 与 y 的两条道路。假设对每个 \(j \in J\)，存在连接 \(\sigma_{j}: I \times I \to X_{j}\) 与 \(\operatorname{pr}_{j} \circ c\) 的严格同伦 \(\operatorname{pr}_{j} \circ d\)。于是，从 \((t,s) \mapsto (\sigma_{j}(t,s))\) 到 X 的映射 \(I \times I\) 是连接 c 与 d 的严格同伦（上引文）。这证明了 \(X\)、\(c\)、\(d\)、\(\varphi_{x,y}\) 是单射。

推论. --- 设 \(x = (x_{j})_{j \in \mathrm{J}}\) 是 X 中的一点。映射

\[
\pi_ {1} (X, x) \to \prod_ {j \in J} \pi_ {1} (X _ {j}, x _ {j})
\]

由映射 \(\pi_1(\mathrm{pr}_j, x_j)\) 诱导，是群同构。

这个同构称为典范同构。下文中，我们常借助此同构将 \(\pi_{1}(\mathrm{X},x)\) 与 \(\prod_{j\in\mathrm{J}}\pi_{1}(\mathrm{X}_{j},x_{j})\) 识别。

注. --- 2) 设 \((\mathrm{X}_{j})_{j\in\mathrm{J}}\) 是一个拓扑空间族。记 X 为积拓扑空间 \(\prod_{j\in J}X_{j}\)，\(x=(x_{j})\) 是 X 中的一点。

对每个 \(i \in \mathrm{J}\)，令 \(u_i: \mathrm{X}_i \to \mathrm{X}\) 为满足下列条件的映射：对 \(z \in \mathrm{X}_i\)，\(\operatorname{pr}_i \circ u_i(z) = z\)，而对所有异于 i 的 \(j \in \mathrm{J}\)，\(i\)、\(\operatorname{pr}_j \circ u_i(z) = x_j\)。映射 \(u_i\) 是连续的，映射 \(\pi_1(u_i, x_i)\) 可与族 \(\pi_1(\mathrm{X}_i, x_i)\) 的积群中因子 \((\pi_1(\mathrm{X}_j, x_j))_{j \in \mathrm{J}}\) 的典范嵌入识别（参见 A, I, p.~45）。

假设集合 J 有限，并且对每个 \(j \in J\)，令 \(\gamma_{j}\) 为 \(\pi_{1}(X_{j}, x_{j})\) 的一个元素。\((\gamma_{j})\) 的元素 \(\pi_{1}(X, x)\) 是圈类 \((u_{j})_{*}(\gamma_{j})\)（\(j \in J\)）的复合，而这些圈类两两可交换。

\begin{enumerate}
\def\labelenumi{\arabic{enumi})}
\setcounter{enumi}{2}
\tightlist
\item
  设 \((\mathrm{X}_j)_{j\in \mathrm{J}}\) 是一个拓扑空间族。记 X 为积拓扑空间 \(\prod_{j\in \mathrm{J}}\mathrm{X}_j\)，\(x = (x_{j})\) 是 X 中的一点。
\end{enumerate}

在集合 \(\pi_1(\mathrm{X},x)\) 和 \(\pi_1(\mathrm{X}_j,x_j)\) 上赋予道路空间 \(\Lambda_x(\mathrm{X})\) 和 \(\Lambda_{x_j}(\mathrm{X}_j)\) 上紧开拓扑的商拓扑。于是同构 \(\pi_1(\mathrm{X},x)\to \prod_{j\in \mathrm{J}}\pi_1(\mathrm{X}_j,x_j)\) 是同胚。它是连续的（III, p.~294, 注 1）。\(\Lambda (\mathrm{X})\) 上的紧开拓扑由形如 \(T(\mathrm{K},\mathrm{U})\) 的子集生成，其中 K 是 I 的紧子集，U 是 X 的开子集。对 \(j\in \mathrm{J}\)，令 \(\mathrm{U}_j\) 为 \(\mathrm{X}_j\) 的开子集，满足 \(\prod_{j\in \mathrm{J}}\mathrm{U}_j\subset \mathrm{U}\)。于是 \((\mathrm{pr}_j)_*(T(\mathrm{K},\mathrm{U}))\) 包含 \(T(\mathrm{K},\mathrm{U}_j)\)。这说明映射 \((\mathrm{pr}_j)_*: \Lambda_x(\mathrm{X})\to \Lambda_{x_j}(\mathrm{X}_j)\) 是开的，映射 \(\pi_1(\mathrm{pr}_j,x_j)\) 也都是开的。由于它们是满射，映射 \(\pi_1(\mathrm{X},x)\to \prod_{j\in \mathrm{J}}\pi_1(\mathrm{X}_j,x_j)\) 是开的（TG, I, p.~34, 命题 8）。它既连续又双射，因而是同胚（TG, I, p.~30, 例 2）。

命题 5. — 设 X 是一个拓扑空间，\((\mathrm{A}_{i})_{i\in\mathrm{I}}\) 是由过滤有序集 I 索引的 X 的子集递增族，并且 X 的每个拟紧子集都包含于某个 \(A_{i}\) 中。典范群胚态射

\[
\rho \colon \varinjlim_ {i \in \mathrm{I}} \varpi (\mathrm{A} _ {i}) \to \varpi (\mathrm{X}),
\]

由 \(A_{i}\) 到 X 的典范嵌入诱导，是同构。若 \(i \leqslant j\)，记 \(\rho_{j,i}\) 为由 \(\varpi(\mathrm{A}_{i}) \to \varpi(\mathrm{A}_{j})\) 到 \(A_{i}\) 的嵌入诱导的群胚态射 \(A_{j}\)。由于 \(\rho_{j,i}\) 取顶点后诱导的映射是嵌入 \(A_{i} \to A_{j}\)，且各 \(A_{i}\) 覆盖 X，\(\rho\) 取顶点后诱导的映射是双射。

设 \(a\) 和 \(b\) 是 X 中的点，\(c\) 是 X 中连接 \(a\) 与 \(b\) 的道路。由于 \(c\) 是紧的，\(\mathbf{I}\) 的像是 X 的拟紧子集（TG, I, p.~62, 定理 2）。因此存在元素 \(i \in \mathrm{I}\)，使得 \(c\) 的像包含于 \(\mathrm{A}_i\) 中。于是，\(\rho\) 取箭集后诱导的映射是满射。

令 \(i \in \mathrm{I}\)，令 \(a\) 和 \(b\) 为 \(\mathrm{A}_i\) 中的点，令 \(c\) 和 \(c'\) 为 \(a\) 中连接 \(b\) 与 \(\mathrm{A}_i\) 的道路；令 \(h\) 为 \(c\) 中连接 \(c'\) 与 \(\mathrm{X}\)、\(\mathbf{I} \times \mathbf{I}\) 的严格同伦。由于 \(h(\mathbf{I} \times \mathbf{I})\) 是紧的，\(\mathrm{X}\) 是 \(i \in \mathrm{I}\) 的拟紧子集（上引文），并且存在元素 \(h\)，使得 h 的像包含于 \(\mathrm{A}_i\)、\(c\) 中。道路 c 与 \(c'\) 在 \(A_{i}\) 中严格同伦；从而，道路类 {[}c{]} 与 \([c']\) 在 \(\varinjlim\varpi(A_{i})\) 中的像相同。因此，\(\rho\) 是单射。

推论. --- 设 a 是 X 中的一点，J 是由满足 \(i \in I\) 的 \(a \in A_{i}\) 组成的集合。典范同态

\[
\varinjlim_ {i \in J} \pi_ {1} (A _ {i}, a) \to \pi_ {1} (X, a)
\]

是双射。

注 4. --- 当 \((\mathrm{A}_{i})_{i\in\mathrm{I}}\) 是由过滤有序集 I 索引的 X 的子集递增族，且各 \(A_{i}\) 的内部覆盖 X 时，本命题及其推论尤其适用。

